\SourceLeaf{23}
\[
\alpha\delta\xi=\delta\alpha.\int\!\int\frac{\mathrm{d}p.\mathrm{d}q.a\sin p}{r^3}\bigl(\frac{(a-x)x}{\alpha\alpha}+\frac{(b-y)y}{\beta\beta}+\frac{(c-z)z}{\gamma\gamma}\bigr)
\]
By equation $[3]$ the right-hand side of this equation is either $=0$ or $=-\frac{4\pi a\delta\alpha}{\alpha\beta\gamma}$, according as $M$ lies outside or inside the body. Thus in the former case we have
\[
[4]\qquad\delta\xi=0
\]
and in the latter case
\[
[5]\qquad\delta\xi=-\frac{4\pi a\delta\alpha}{\alpha\alpha\beta\gamma}
\]
\par
Equation $[4]$ shows at once that $\xi$ is constant, or that the attraction $X$ is proportional to the mass, for all ellipsoids in which $\alpha\alpha-\beta\beta$ and $\alpha\alpha-\gamma\gamma$ are constant quantities, that is, whose three principal sections are ellipses drawn from the same foci, as long as the attracted point lies outside the spheroid. Since this conclusion is rigorously true however close the surface of the spheroid comes to the attracted point, it must also be possible to extend it to the spheroid itself whose surface passes through the attracted point.

The problem of determining the attraction of a spheroid on any exterior point is thus reduced to two other problems. The first is to determine the dimensions of another spheroid drawn from the same foci as the given one and passing through the attracted point. The second is to determine the attraction of a spheroid on a point on its surface. The first problem depends on solving a cubic equation. It is easy to show that this equation always has exactly one real root, and there seems to be no need to dwell on it here. To solve the second problem, however, let us consider the second case, in which the attracted point lies inside the body. Since $\beta\beta=\alpha\alpha+BB-AA$, $\gamma\gamma=\alpha\alpha+CC-AA$, we will substitute these values into equation $5$ and at the same time set $\frac{A}{\alpha}=t$. This yields
\[
\delta\xi=\frac{4a\pi tt\delta t}{A^3\BareRoot{((1-(1-\frac{BB}{AA})tt)(1-(1-\frac{CC}{AA})tt)}}
\]
or, returning to the differential sign $\mathrm{d}$ and integrating,
\[
\xi=\frac{4a\pi}{A^3}\int\frac{tt\mathrm{d}t}{\BareRoot{((1-(1-\frac{BB}{AA})tt)(1-(1-\frac{CC}{AA})tt)}}
\]
The integral is to be taken so that it vanishes for $t=0$, and then, for the particular spheroid whose semiaxes are $A,B,C$, it is to be extended as far as $t=1$. We thus have
